\section{Related Work}
\subsection{Adversarial Attacks and Training}
FGSM follows the sign of the input loss gradient in one step
\cite{goodfellow2015explaining}. PGD iterates gradient steps with projection
onto a prescribed threat set and is central to the robust-optimization view of
adversarial training \cite{madry2018towards}. TRADES separates natural error
from a boundary-related robustness term and makes the clean--robust trade-off
explicit \cite{zhang2019trades}. Although these methods are effective, their
multi-step inner loops increase training cost. Free adversarial training reuses
gradients across repeated mini-batches \cite{shafahi2019adversarial}, whereas
fast training uses a single randomized FGSM step \cite{wong2020fast}.

Robustness evaluation itself must be sufficiently strong. PGD results can be
sensitive to iterations, restarts, and objectives. AutoAttack combines
complementary parameter-free attacks to reduce evaluator bias and expose
gradient masking \cite{croce2020reliable}. In the controlled study below, we
use progressively stronger full-test PGD variants and multiple restarts rather
than treating FGSM accuracy as a robustness estimate.

\subsection{Catastrophic Overfitting}
Wong \emph{et al.} identified CO as a failure mode of fast adversarial
training \cite{wong2020fast}. Andriushchenko and Flammarion connected it to
local nonlinearity and proposed GradAlign, which regularizes gradient alignment
inside the perturbation set \cite{andriushchenko2020understanding}. Li
\emph{et al.} showed that recovery from weak-attack overfitting is important
and proposed FastAdv+ and FastAdvW, which use periodic PGD validation to decide
when to introduce PGD updates \cite{li2020towards}. Our response mechanism is
related in spirit, but the detector differs operationally: the endpoint input
gradient is obtained from the same parameter-backward pass, without an
additional forward or backward pass. The monitor does not execute a validation
attack. We also separate prediction from mitigation by freezing the detector
before held-out evaluation.

\subsection{Empirical and Certified Robustness}
Empirical defenses are evaluated against known attacks but do not guarantee
correctness for every allowed perturbation. Certified methods instead compute a
provable lower bound, for example through convex outer approximations
\cite{wong2018provable} or randomized smoothing under an $\ell_2$ threat model
\cite{cohen2019certified}. This paper targets efficient empirical training in
the $\linf$ setting. The distinction is important: the attack results reported
below measure resistance to a strong test suite, not certified radii.
